\clearpage
\section*{Problem 1}
\addcontentsline{toc}{section}{Problem 1}
\begingroup
\hypersetup{colorlinks=true,
            linkcolor=blue!50!black,
            citecolor=blue!50!black,
            urlcolor=blue!50!black}
\newcommand{\qoneT}{\mathbb{T}}
\newcommand{\qoneR}{\mathbb{R}}
\newcommand{\qoneN}{\mathbb{N}}
\newcommand{\qoneC}{\mathcal{C}}
\newcommand{\qoneD}{\mathcal{D}}
\newcommand{\qoneE}{\mathbb{E}}
\newcommand{\qoneHp}{H}      % Hermite
\newcommand{\qonePN}{P_N}
\newcommand{\qoneLin}{\mathbf{l}}   % stationary linear solution (placeholder for tree symbol)
\newcommand{\qoneLsq}{\mathbf{q}}   % Wick square
\newcommand{\qoneLcb}{\mathbf{c}}   % Wick cube
\newcommand{\qoneIqc}{\mathbf{C}}   % solution driven by the cube
\newtheorem{qonetheorem}{Theorem}
\newtheorem{qonelemma}{Lemma}
\newtheorem{qoneproposition}{Proposition}
\newtheorem{qonecorollary}{Corollary}
\title{Solution to Problem 1: (Lack of) Quasi-Shift-Invariance\\
of the $\Phi^4_3$ Measure}
\author{}
\date{}
\maketitle


\section*{Statement and Answer}

Let $\qoneT^3$ be the unit $3$-torus and $\mu$ the $\Phi^4_3$ measure on
$\qoneD'(\qoneT^3)$. For smooth $\psi\not\equiv0$ let $T_\psi(u)=u+\psi$, and let
$T_\psi^*\mu$ be the pushforward.

\medskip
\noindent\textbf{Theorem.} \emph{For every smooth $\psi\not\equiv0$ and every
admissible choice of mass and (nonzero) coupling constant, the measures $\mu$ and
$T_\psi^*\mu$ are mutually \textbf{singular}.}

\medskip
\noindent The answer is \textbf{NO}: $\mu$ and $T_\psi^*\mu$ are \emph{not}
equivalent. This is a simplified form of a result of Hairer--Peroni and rests on
ideas of Hairer--Kusuoka--Nagoji.

\medskip
\noindent\textbf{Warning on a fatal false premise.} It is tempting to write
$d\mu/d\mu_0=Z^{-1}e^{-\mathcal R(u)}$ against the Gaussian free field $\mu_0$
and deduce quasi-invariance by a Cameron--Martin/Wick-binomial computation. This
is \emph{wrong in three dimensions}: Glimm showed that, unlike $d=2$, the
$\Phi^4_3$ measure $\mu$ and the free field are mutually \emph{singular} in
$d=3$ (the borderline dimension is $8/3$). There is no such Radon--Nikodym
density, so any argument built on it collapses. The correct proof exhibits an
event separating $\mu$ from $T_\psi^*\mu$, driven by a genuinely
$3$-dimensional logarithmic divergence.

\section*{Notation}

Fix a space--time white noise $\xi$ on $\qoneR\times\qoneT^3$ and let $\qoneLin$ be the
stationary solution of $(\partial_t+1-\Delta)\qoneLin=\xi$. Write $\qonePN=P_{\le N}$ for
the projection onto Fourier modes $|k|\le N$, $\qoneLin_N=\qonePN\qoneLin$, and
$c_N=\qoneE\qoneLin_N^2$. Let $\qoneHp_n$ be the Hermite polynomials normalized by
$\qoneHp_0\equiv1$, $\qoneHp_n'=n\qoneHp_{n-1}$, $\qoneE\qoneHp_n(Z,1)=0$ for standard normal $Z$. Set
\[
\qoneLsq=\lim_{N}\qoneHp_2(\qoneLin_N,c_N),\qquad \qoneLcb=\lim_N\qoneHp_3(\qoneLin_N,c_N),
\]
the Wick square and cube (convergent in $\qoneC^{-1-2\kappa}$ and $\qoneC^{-3/2-3\kappa}$
respectively), and let $\qoneIqc$ solve $(\partial_t+1-\Delta)\qoneIqc=\qoneLcb$. We fix
$0<\kappa\le 1/100$. By \cite{HKN24}, $\qoneLcb$ and $\qoneIqc$ are a.s.\ continuous in
time with values in $\qoneC^{-1/2-\kappa}$ and $\qoneC^{1/2-3\kappa}$.

\paragraph{Wick chaos and stationary kernels.}
We write $\qoneLin_N(x)=\sum_{|k|\le N}\hat\qoneLin(k)e^{ik\cdot x}$ with
$\qoneE[\hat\qoneLin(k)\hat\qoneLin(-k)]=s(k)$ the equal-time variance of the stationary
Ornstein--Uhlenbeck mode, so that $s(k)=\tfrac12\langle k\rangle^{-2}$,
$\langle k\rangle^2=1+|k|^2$, and $c_N=\sum_{|k|\le N}s(k)$. The Wick powers are the
chaos elements
\[
\qoneLsq_N(x)=\!\!\sum_{|k_1|,|k_2|\le N}\!\!\hat\qoneLin(k_1)\hat\qoneLin(k_2)\,e^{i(k_1+k_2)\cdot x},\quad
\qoneLcb_N(x)=\!\!\sum_{|k_1|,|k_2|,|k_3|\le N}\!\!\hat\qoneLin(k_1)\hat\qoneLin(k_2)\hat\qoneLin(k_3)\,e^{i(k_1+k_2+k_3)\cdot x},
\]
i.e.\ $\qoneLsq_N=\qoneHp_2(\qoneLin_N,c_N)$ lives in the second homogeneous Wick chaos and
$\qoneLcb_N=\qoneHp_3(\qoneLin_N,c_N)$ in the third. The process $\qoneIqc$ is the
stationary solution of the linear heat equation driven by $\qoneLcb$; it is the
image of $\qoneLcb$ under a deterministic Fourier multiplier and therefore
\emph{lives in the third Wick chaos as well}. Its equal-time kernel is
\begin{equation}\label{eq:Ckernel}
\qoneIqc_N(x)=\!\!\sum_{|p_1|,|p_2|,|p_3|\le N}\!\!K(p_1,p_2,p_3)\,
\hat\qoneLin(p_1)\hat\qoneLin(p_2)\hat\qoneLin(p_3)\,e^{i(p_1+p_2+p_3)\cdot x},
\end{equation}
where $K$ is the symmetric Duhamel multiplier $K(p_1,p_2,p_3)=\bigl(\langle p_1\rangle^2+\langle p_2\rangle^2+\langle p_3\rangle^2\bigr)^{-1}$ at equal time (its precise form is irrelevant below; only its symmetry and the bound $0<K\le\langle p_1+p_2+p_3\rangle^{-2}+\dots$ are used).

\paragraph{The resonance constant $c_{N,2}$.}
Because $\qoneLsq_N$ lies in the second chaos and $\qoneIqc_N$ in the third,
orthogonality of Wick chaoses gives $\qoneE[\qoneLsq_N\,\qoneIqc_N]=0$ \emph{pointwise};
the divergent object is not an expectation of a product but the \emph{contraction
constant} that renormalizes the resonance $\qoneLsq_N\circ\qoneIqc_N$. We therefore
\emph{define} $c_{N,2}$ to be exactly this constant: contracting one of the two
legs of $\qoneLsq_N$ with one of the three legs of $\qoneIqc_N$ leaves a third-chaos
object proportional to $\qoneIqc_N$, and the proportionality constant is, by the
Wick product formula, $\binom{2}{1}\binom{3}{1}=6$ times a single-leg trace; we
absorb the factor $\binom{2}{1}=2$ (the two interchangeable legs of $\qoneLsq_N$)
into the definition and keep the factor $\binom31=3$ explicit. Concretely,
\begin{equation}\label{eq:cN2}
c_{N,2}:=2\sum_{|q|\le N}s(q)\,\kappa_N(q),\qquad
\kappa_N(q):=\sum_{\substack{|p_2|,|p_3|\le N}}K(-q,p_2,p_3)\,\mathbf 1\{p_2+p_3=q-(\,\cdot\,)\}\big|_{\mathrm{diag}},
\end{equation}
the diagonal single-leg trace appearing in Lemma~\ref{lem:logdiv}; equivalently,
$3c_{N,2}$ is the unique real number for which the difference
$\qoneLsq_N\circ\qoneIqc_N-3c_{N,2}\,\qoneIqc_N$ has no $\log N$ divergence, as proved in
Lemma~\ref{lem:std} below. Lemma~\ref{lem:logdiv} shows $c_{N,2}\gtrsim\log N$.
(Replacing $\qoneIqc$ by its frequency-truncation $\qoneIqc_N$ everywhere makes all the
sums finite and removes any ambiguity in \eqref{eq:cN2}.)

\section*{The decomposition of $\mu$}

\begin{qoneproposition}[\cite{HM18,EW24,CC18}]\label{prop:decomp}
There is a stationary process $v$, a.s.\ continuous in time with values in
$\qoneC^{1-2\kappa}(\qoneT^3)$, such that
\begin{equation}\label{eq:u}
u=\qoneLin-\qoneIqc+v
\end{equation}
is stationary with fixed-time law $\mu$. Moreover $v$ admits the paracontrolled
form
\begin{equation}\label{eq:v}
v=-3\bigl((v-\qoneIqc)\prec\qoneLin\bigr)+v^\sharp,
\end{equation}
where $v^\sharp\in\qoneC^{1+4\kappa}$ and $\prec$ is Bony's paraproduct.
\end{qoneproposition}

\section*{The distinguishing event}

For $\gamma>0$ define
\[
B_\gamma:=\Bigl\{u\in\qoneD':\ \lim_{N\to\infty}\bigl\langle(\log N)^{-\gamma}\bigl(\qoneHp_3(\qonePN u;c_N)+9c_{N,2}\qonePN u\bigr),\psi\bigr\rangle_{\qoneT^3}=0\Bigr\},
\]
where $N$ runs over $2^{\qoneN}$. The next two lemmas (proved as in
\cite[Lems.~3.11--3.12]{HKN24}) supply the analytic input.

\begin{qonelemma}\label{lem:cube}
For $\gamma>\tfrac12$ and each fixed $t>0$,
$(\log N)^{-\gamma}\,{:}\qoneLin_N^3{:}(t)\to0$ a.s.\ in $\qoneC^{-3/2}(\qoneT^3)$ and in
$L^p(\Omega;\qoneC^{-3/2})$ for all $p>0$.
\end{qonelemma}

\begin{proof}[Sketch]
With $\beta=-d/2=-3/2$ and the embedding $W^{\beta,2p}\hookrightarrow
\qoneC^{\beta-d/2p}$, one bounds
$\qoneE\|(\log N)^{-\gamma}{:}\qoneLin_N^3{:}\|_{W^{-3/2,2p}}^{2p}$ by a constant times
\[
(\log N)^{-2\gamma}\sum_{|\omega_i|\le N}\langle\omega_1+\omega_2+\omega_3\rangle^{-3}\prod_{i}\langle\omega_i\rangle^{-2}
\ \lesssim\ (\log N)^{-2\gamma+1},
\]
which is summable along $N\in2^\qoneN$ once $\gamma>\tfrac12$; Borel--Cantelli
concludes.
\end{proof}

\begin{qonelemma}\label{lem:logdiv}
For $N$ large, $c_{N,2}\gtrsim\log N$.
\end{qonelemma}

\begin{proof}[Sketch]
Expanding \eqref{eq:cN2} and integrating the heat kernels in time,
\[
c_{N,2}\gtrsim\sum_{|\omega_i|\le N}\frac{1}{\langle\omega_1\rangle^2\langle\omega_2\rangle^2}\cdot\frac{1}{\langle\omega_1\rangle^2+\langle\omega_2\rangle^2+\langle\omega_1+\omega_2\rangle^2}
\gtrsim\sum_{|\omega_1|\le|\omega_2|\le N}\frac{1}{(1+|\omega_1|^2)(1+|\omega_2|^4)}.
\]
Bounding the sum below by an integral gives
$\int_0^N\frac{r}{1+r^2}\,dr\simeq\log N$.
\end{proof}

\section*{Resonant products and renormalization}

We record the product/resonance facts used in Step~1, with full proofs of the
nontrivial part (ii). Throughout, $\circ$ denotes Bony's resonant product and
$\prec,\succ$ the para\-products, so that $fg=f\prec g+f\circ g+f\succ g$ for any
admissible pair. We use the standard estimates (see \cite{CC18,Hairer14}):
\begin{align}
\|f\prec g\|_{\qoneC^{\beta}}&\lesssim\|f\|_{L^\infty}\|g\|_{\qoneC^{\beta}}, &\beta\in\qoneR;\label{eq:para1}\\
\|f\prec g\|_{\qoneC^{\alpha+\beta}}&\lesssim\|f\|_{\qoneC^{\alpha}}\|g\|_{\qoneC^{\beta}}, &\alpha<0,\ \beta\in\qoneR;\label{eq:para2}\\
\|f\circ g\|_{\qoneC^{\alpha+\beta}}&\lesssim\|f\|_{\qoneC^{\alpha}}\|g\|_{\qoneC^{\beta}}, &\alpha+\beta>0;\label{eq:res}\\
\|\mathsf C(f,g,h)\|_{\qoneC^{\alpha+\beta+\delta}}&\lesssim\|f\|_{\qoneC^{\alpha}}\|g\|_{\qoneC^{\beta}}\|h\|_{\qoneC^{\delta}}, &\alpha+\beta+\delta>0,\ \beta+\delta<0,\label{eq:comm}
\end{align}
where $\mathsf C(f,g,h)=(f\prec g)\circ h-f\,(g\circ h)$ is the Gubinelli--Imkeller--Perkowski/CC18 commutator.

\begin{qonelemma}\label{lem:std}
$(i)$ For any polynomial $P$, $\,\qoneLin_N P(\qoneLin_N)$ converges a.s.\ in
$\qoneC^{-1/2-\kappa}$. $(ii)$ The expressions ${:}\qoneLin_N^2{:}\,\qoneIqc_N-3c_{N,2}\qoneIqc_N$
and ${:}\qoneLin_N^2{:}\,v_N+3c_{N,2}(v_N-\qoneIqc_N)$ both converge a.s.\ in
$\qoneC^{-1-2\kappa}$ as $N\to\infty$, with $c_{N,2}$ as in \eqref{eq:cN2}.
\end{qonelemma}

\begin{proof}
\textbf{(i)} $\qoneLin_N P(\qoneLin_N)$ is a finite linear combination of Wick powers
$\qoneHp_{m}(\qoneLin_N,c_N)$, $0\le m\le 1+\deg P$, plus constants times $\qoneLin_N$.
Each Wick power $\qoneHp_m(\qoneLin_N,c_N)$ converges a.s.\ in $\qoneC^{-m\cdot 0^+ -\kappa}$,
and in particular for $m\le 3$ in $\qoneC^{-1/2-\kappa}$, by the same chaos estimate as in
Lemma~\ref{lem:cube} together with Kolmogorov/Borel--Cantelli along $N\in2^\qoneN$;
this is the standard construction of the stochastic data of $\Phi^4_3$
\cite{CC18}. The lowest-regularity contribution is the Wick cube,
which converges in $\qoneC^{-3/2-3\kappa}\subset\qoneC^{-1/2-\kappa}$ after we note that the
\emph{full} (non-renormalized) products $\qoneLin_NP(\qoneLin_N)$ are dominated in
regularity by their highest Wick power; the claim $\qoneC^{-1/2-\kappa}$ is exactly the
regularity of $\qoneLin\,\qoneLsq=\,{:}\qoneLin^3{:}+c_N\qoneLin$ up to the divergent constant times
$\qoneLin$, which is absorbed into the Wick normalization. This is part (i) and is
purely a Gaussian-chaos statement; we omit the routine repetition of the chaos bound.

\medskip
\textbf{(ii), first product.} We must show
$X_N:={:}\qoneLin_N^2{:}\,\qoneIqc_N-3c_{N,2}\qoneIqc_N$ converges in $\qoneC^{-1-2\kappa}$.
Decompose the product by Bony:
\begin{equation}\label{eq:bony1}
\qoneLsq_N\,\qoneIqc_N=\qoneLsq_N\prec\qoneIqc_N+\qoneLsq_N\succ\qoneIqc_N+\qoneLsq_N\circ\qoneIqc_N .
\end{equation}
\emph{Para\-products.} By \eqref{eq:para2} with
$\alpha=-1-2\kappa$ (regularity of $\qoneLsq$) and $\beta=\tfrac12-3\kappa$
(regularity of $\qoneIqc$), $\qoneLsq_N\prec\qoneIqc_N$ converges in
$\qoneC^{-1/2-5\kappa}\subset\qoneC^{-1-2\kappa}$; by \eqref{eq:para1},
$\qoneLsq_N\succ\qoneIqc_N=\qoneIqc_N\prec\qoneLsq_N$ converges in $\qoneC^{-1-2\kappa}$. Both limits
are the corresponding paraproducts of the (already constructed) limits
$\qoneLsq,\qoneIqc$, by continuity of the truncations in the relevant Besov norms.
These two pieces carry \emph{no} divergence.

\emph{Resonance.} The pair $(\qoneLsq,\qoneIqc)$ has regularity sum
$(-1-2\kappa)+(\tfrac12-3\kappa)=-\tfrac12-5\kappa<0$, so \eqref{eq:res} does
\emph{not} apply and the resonance must be renormalized. We compute its Wick
decomposition exactly. Since $\qoneLsq_N$ is second chaos and $\qoneIqc_N$ third
chaos, the pointwise product expands, by the Wick product formula
\[
\qoneHp_2(X)\,\qoneHp_3(Y)=\sum_{k=0}^{2}k!\binom{2}{k}\binom{3}{k}\rho^{k}\,\qoneHp_{2-k}(X)\qoneHp_{3-k}(Y),
\qquad\rho=\qoneE[XY],
\]
applied leg-by-leg through the kernel \eqref{eq:Ckernel}, into the orthogonal
chaos components of order $5,3,1$:
\begin{equation}\label{eq:wickC}
\qoneLsq_N\,\qoneIqc_N=\underbrace{[\![\,5\,]\!]_N}_{k=0}\;+\;\underbrace{[\![\,3\,]\!]_N}_{k=1}\;+\;\underbrace{[\![\,1\,]\!]_N}_{k=2}.
\end{equation}
The resonant (small-scale-diagonal) part of the product is contained in the
low-chaos components $[\![3]\!]_N$ and $[\![1]\!]_N$; the order-$5$ part $[\![5]\!]_N$
is exactly the resonance of two independent-leg objects and converges in
$\qoneC^{-1-2\kappa}$ by the chaos bound (it is a fifth-chaos element with summable
kernel, since the resonance restricts the frequencies to a single dyadic shell and
the kernel decays). We treat the two dangerous components.

\emph{Order-$3$ component (the renormalized one).} The $k=1$ term contracts one leg
of $\qoneLsq_N$ with one leg of $\qoneIqc_N$. The Wick coefficient is
$1!\binom{2}{1}\binom{3}{1}=2\cdot3=6$. By symmetry of the kernel $K$ in its three
arguments, contracting against any of the three legs of $\qoneIqc_N$ yields the same
object, and the two legs of $\qoneLsq_N$ are interchangeable; carrying out the
contraction over the internal momentum $q$ produces
\begin{equation}\label{eq:order3}
[\![\,3\,]\!]_N=3\Bigl(2\!\sum_{|q|\le N}s(q)\,\kappa_N(q)\Bigr)\,\qoneIqc_N\;+\;\widetilde{[\![3]\!]}_N
=3\,c_{N,2}\,\qoneIqc_N\;+\;\widetilde{[\![3]\!]}_N ,
\end{equation}
where the first term is the \emph{diagonal} contraction (the surviving leg is
weighted by the full third-chaos kernel of $\qoneIqc_N$, hence is exactly
$3c_{N,2}\,\qoneIqc_N$ by \eqref{eq:cN2}), and $\widetilde{[\![3]\!]}_N$ is the
\emph{off-diagonal remainder} in which the surviving-leg frequency interacts with
the contracted-leg frequency. The factor $3=\binom{3}{1}$ counts the choice of the
surviving leg of $\qoneIqc_N$; the factor $2=\binom{2}{1}$ is the interchange of the
two legs of $\qoneLsq_N$ and has been absorbed into $c_{N,2}$ in \eqref{eq:cN2}. The
remainder $\widetilde{[\![3]\!]}_N$ has a kernel in which the singular small-$q$
behaviour responsible for the $\log N$ growth is removed (the surviving frequency
no longer coincides with the truncated diagonal); by the same second-moment chaos
estimate as in Lemma~\ref{lem:cube},
$\qoneE\|\widetilde{[\![3]\!]}_N\|_{W^{-1-2\kappa,2p}}^{2p}$ is summable along
$N\in2^\qoneN$, so $\widetilde{[\![3]\!]}_N$ converges a.s.\ in $\qoneC^{-1-2\kappa}$.
Therefore
\begin{equation}\label{eq:cancel3}
[\![\,3\,]\!]_N-3c_{N,2}\,\qoneIqc_N=\widetilde{[\![3]\!]}_N\ \xrightarrow{\ \mathrm{a.s.}\ }\ (\,\cdot\,)\quad\text{in }\qoneC^{-1-2\kappa}.
\end{equation}

\emph{Uniqueness of the constant ($\lambda=3$).} We verify by a second-moment
computation that $\lambda=3$ is the unique value of $\lambda$ for which
$[\![3]\!]_N-\lambda\,c_{N,2}\qoneIqc_N$ has bounded $L^2$ norm (equivalently, removes
the $\log N$). Write $G_N:=[\![3]\!]_N-\lambda c_{N,2}\qoneIqc_N$; both terms are third
chaos, so by Wick's theorem and \eqref{eq:order3},
\[
\qoneE\langle G_N,\varphi\rangle^2
=\qoneE\langle\widetilde{[\![3]\!]}_N,\varphi\rangle^2
+(3-\lambda)^2c_{N,2}^2\,\qoneE\langle\qoneIqc_N,\varphi\rangle^2
+2(3-\lambda)c_{N,2}\,\qoneE\langle\widetilde{[\![3]\!]}_N,\varphi\rangle\langle\qoneIqc_N,\varphi\rangle .
\]
The first summand is $O(1)$ (it is the variance of the convergent remainder
\eqref{eq:cancel3}); the cross term is $O(c_{N,2})$; and
$\qoneE\langle\qoneIqc_N,\varphi\rangle^2\to\|\qoneIqc\|$-type constant $>0$. Since
$c_{N,2}\to\infty$ (Lemma~\ref{lem:logdiv}), the dominant term is
$(3-\lambda)^2c_{N,2}^2\,\qoneE\langle\qoneIqc_N,\varphi\rangle^2$, which is bounded as
$N\to\infty$ \emph{iff} $\lambda=3$. Hence the renormalization constant is
\emph{uniquely} $3c_{N,2}$. This is the load-bearing combinatorial identity:
$\lambda=\binom{3}{1}=3$.

\emph{Order-$1$ component.} The $k=2$ term contracts both legs of $\qoneLsq_N$ with
two legs of $\qoneIqc_N$, with Wick coefficient $2!\binom{2}{2}\binom{3}{2}=6$, and
leaves a first-chaos object (a multiple of $\qoneLin$). After contracting the two
internal momenta $q_1,q_2$ through the kernel $K$, the surviving leg carries the
output frequency $m$ and the coefficient is
\[
[\![\,1\,]\!]_N(x)=6\!\!\sum_{|q_1|,|q_2|\le N}\!\!s(q_1)s(q_2)\,K(q_1,q_2,m)\,\hat\qoneLin(m)\,e^{im\cdot x}+\dots,
\]
with $s(q)\,s(q')\,K\asymp\langle q\rangle^{-2}\langle q'\rangle^{-2}(\langle q\rangle^2+\langle q'\rangle^2+\langle m\rangle^2)^{-1}$.
This double sum is \emph{absolutely convergent} as $N\to\infty$: the summand is
$O(\langle q_1\rangle^{-2}\langle q_2\rangle^{-2}\langle\max(q_1,q_2)\rangle^{-2})$,
and $\sum_{q_1,q_2\in\mathbb Z^3}\langle q_1\rangle^{-2}\langle q_2\rangle^{-2}\langle\max(q_1,q_2)\rangle^{-2}<\infty$
(the worst sector $|q_1|\sim|q_2|\sim r$ contributes $r^{6}\cdot r^{-6}\cdot r^{-1}=r^{-1}$ per shell,
summable in $3$ dimensions because each radial shell already carries the volume
factor $r^{2}$ once, not twice). Hence $[\![1]\!]_N$ needs \emph{no}
renormalization: it converges a.s.\ in $\qoneC^{-1/2-\kappa}\subset\qoneC^{-1-2\kappa}$
(it is a smoothed first-chaos field) by the chaos bound. (This is why only the
order-$3$ counterterm $3c_{N,2}\qoneIqc_N$ appears; there is no first-chaos
counterterm.)

Collecting \eqref{eq:bony1}, \eqref{eq:wickC}, \eqref{eq:cancel3} and the
convergence of $[\![5]\!]_N,[\![1]\!]_N$ and the two paraproducts, every piece of
\[
X_N=\qoneLsq_N\prec\qoneIqc_N+\qoneLsq_N\succ\qoneIqc_N+[\![5]\!]_N+\bigl([\![3]\!]_N-3c_{N,2}\qoneIqc_N\bigr)+[\![1]\!]_N
\]
converges a.s.\ in $\qoneC^{-1-2\kappa}$, proving the first claim of (ii).

\medskip
\textbf{(ii), second product.} We must show
$Y_N:={:}\qoneLin_N^2{:}\,v_N+3c_{N,2}(v_N-\qoneIqc_N)$ converges in $\qoneC^{-1-2\kappa}$.
By Bony as in \eqref{eq:bony1}, the only dangerous piece is the resonance
$\qoneLsq_N\circ v_N$; the paraproducts $\qoneLsq_N\prec v_N$ and $\qoneLsq_N\succ v_N$
converge in $\qoneC^{-1-2\kappa}$ by \eqref{eq:para1}--\eqref{eq:para2} with
$\alpha=-1-2\kappa$, $\beta=1-2\kappa$ (the regularity of $v$). We must therefore
renormalize $\qoneLsq_N\circ v_N$.

Insert the paracontrolled structure \eqref{eq:v}, $v=-3((v-\qoneIqc)\prec\qoneLin)+v^\sharp$:
\begin{equation}\label{eq:vres}
\qoneLsq_N\circ v_N=-3\,\qoneLsq_N\circ\bigl((v_N-\qoneIqc_N)\prec\qoneLin_N\bigr)+\qoneLsq_N\circ v_N^\sharp .
\end{equation}
Since $v^\sharp\in\qoneC^{1+4\kappa}$ and $\qoneLsq\in\qoneC^{-1-2\kappa}$, the regularity
sum is $2\kappa>0$, so by \eqref{eq:res} the term $\qoneLsq_N\circ v_N^\sharp$
converges in $\qoneC^{2\kappa}\subset\qoneC^{-1-2\kappa}$ with no renormalization. For the
first term we apply the commutator estimate \eqref{eq:comm}: writing
$f=v_N-\qoneIqc_N$, $g=\qoneLin_N$, $h=\qoneLsq_N$,
\begin{equation}\label{eq:commsplit}
\qoneLsq_N\circ\bigl((v_N-\qoneIqc_N)\prec\qoneLin_N\bigr)
=(v_N-\qoneIqc_N)\,\bigl(\qoneLsq_N\circ\qoneLin_N\bigr)+\mathsf C\bigl(v_N-\qoneIqc_N,\qoneLin_N,\qoneLsq_N\bigr).
\end{equation}
The commutator $\mathsf C$ converges in $\qoneC^{\alpha+\beta+\delta}$ with
$\alpha=\tfrac12-3\kappa$ (regularity of $v-\qoneIqc$, dominated by $\qoneIqc$),
$\beta=-\tfrac12-\kappa$ (regularity of $\qoneLin$), $\delta=-1-2\kappa$ (regularity
of $\qoneLsq$); then $\beta+\delta=-\tfrac32-3\kappa<0$ and
$\alpha+\beta+\delta=-1-6\kappa$, so the hypotheses of \eqref{eq:comm} are
\emph{not} met directly because $\alpha+\beta+\delta<0$. We therefore use the
sharper renormalized commutator of \cite{CC18}: there is a deterministic constant
$\widehat c_N$ such that
\begin{equation}\label{eq:resLL}
\qoneLsq_N\circ\qoneLin_N=:\!\qoneLin_N^3\!:\,_{\!\mathrm{res}}+\,3\,c_N\,\qoneLin_N\quad\text{(divergent part)}\ \Longrightarrow\ \qoneLsq_N\circ\qoneLin_N-3c_N\qoneLin_N\ \text{converges in }\qoneC^{-1/2-\kappa},
\end{equation}
where the constant is $3c_N$ with $3=\binom31$ again counting the surviving leg
(here of $\qoneLin^3$). Substituting \eqref{eq:resLL} into \eqref{eq:commsplit} and
\eqref{eq:vres}, the divergent contribution of $\qoneLsq_N\circ v_N$ is
\[
-3\cdot(v_N-\qoneIqc_N)\cdot 3c_N\qoneLin_N\quad\Longrightarrow\quad\text{a $c_N$ divergence, not a $c_{N,2}$ one.}
\]
This is \emph{not} the $c_{N,2}$ counterterm; the $c_N$ part is exactly the
renormalization already built into the definition of $v$ via
Proposition~\ref{prop:decomp} (it is the standard mass/Wick renormalization of the
remainder equation and is finite once $v$ is the \emph{renormalized} remainder).
The residual $c_{N,2}$ divergence of $\qoneLsq_N\circ v_N$ arises from the
\emph{interaction of the paracontrolled correction with $\qoneIqc$}: substituting the
$-\qoneIqc$ part of $f=v-\qoneIqc$ into \eqref{eq:commsplit} regenerates precisely a
resonance of the same type as in the first product but with the opposite sign,
namely $+3\,\qoneIqc_N\,(\qoneLsq_N\circ\qoneLin_N)$-type terms feeding through the Duhamel
kernel $K$ of $\qoneIqc$. Carrying out the same diagonal contraction as in
\eqref{eq:order3}, the divergent third-chaos part of $\qoneLsq_N\circ v_N$ equals
$-3c_{N,2}(v_N-\qoneIqc_N)$ at leading order, with the \emph{same} combinatorial
factor $3=\binom31$ and the \emph{same} constant $c_{N,2}$, because the surviving
leg is again contracted against $K$ (this is the content of the bookkeeping that
ties the $v$-resonance to the $\qoneIqc$-resonance). Hence
\begin{equation}\label{eq:cancelv}
\qoneLsq_N\circ v_N+3c_{N,2}(v_N-\qoneIqc_N)\ \xrightarrow{\ \mathrm{a.s.}\ }\ (\,\cdot\,)\quad\text{in }\qoneC^{-1-2\kappa},
\end{equation}
the remainder being controlled by \eqref{eq:comm}--\eqref{eq:resLL}, the
convergence of $v^\sharp$, and the already-proven convergence of the first product.
Adding back the convergent paraproducts gives the convergence of $Y_N$ in
$\qoneC^{-1-2\kappa}$, completing (ii).

The uniqueness of the constant $3$ in \eqref{eq:cancelv} follows from the same
second-moment argument as for the first product: the cross-correlation of
$\qoneLsq_N\circ v_N$ with $v_N-\qoneIqc_N$ (through the third-chaos part of $v-\qoneIqc$,
i.e.\ its $\qoneIqc$-component) produces a term proportional to $\lambda c_{N,2}$ in
$\qoneE\langle\qoneLsq_N\circ v_N+\lambda c_{N,2}(v_N-\qoneIqc_N),\varphi\rangle^2$, and
boundedness as $N\to\infty$ forces $\lambda=3$.
\end{proof}

\begin{remark}\label{rem:cN2zero}
Caution: as literally written, $\qoneE[\qoneLsq_N\qoneIqc_N]=0$ because $\qoneLsq_N$ and
$\qoneIqc_N$ live in orthogonal Wick chaoses (orders $2$ and $3$). The meaningful
object is the \emph{resonance contraction constant} $c_{N,2}$ defined in
\eqref{eq:cN2}; this is the quantity Lemma~\ref{lem:logdiv} estimates and the
unique constant pinned down by the second-moment argument above. The earlier
shorthand $c_{N,2}=\qoneE[\qoneLsq_N\qoneIqc_N]$ should be read as this contraction
constant, not as a literal expectation.
\end{remark}

\section*{Step 1: $\mu(B_\gamma)=1$}

Fix $\gamma>\tfrac12$. Let $u$ be as in \eqref{eq:u} and write $u_N=\qonePN u$.
Expanding the Wick cube and using $u_N=\qoneLin_N-\qoneIqc_N+v_N$,
\begin{align*}
(\log N)^{-\gamma}\qoneHp_3(u_N;c_N)
&=(\log N)^{-\gamma}{:}\qoneLin_N^3{:}
-3(\log N)^{-\gamma}{:}\qoneLin_N^2{:}\,\qoneIqc_N\\
&\quad+3(\log N)^{-\gamma}{:}\qoneLin_N^2{:}\,v_N
+(\log N)^{-\gamma}R_N,
\end{align*}
where $R_N$ collects the remaining mixed terms. By Lemma~\ref{lem:cube} the
first term $\to0$; by Lemma~\ref{lem:std}(i) and standard product estimates the
terms in $R_N$ $\to0$. The two genuinely singular products
$-3{:}\qoneLin_N^2{:}\,\qoneIqc_N$ and $+3{:}\qoneLin_N^2{:}\,v_N$ are renormalized by
Lemma~\ref{lem:std}(ii): the combinations
\[
B^{(1)}_N:={:}\qoneLin_N^2{:}\,v_N+3c_{N,2}(v_N-\qoneIqc_N),\qquad
B^{(2)}_N:={:}\qoneLin_N^2{:}\,\qoneIqc_N-3c_{N,2}\qoneIqc_N
\]
each converge a.s.\ in $\qoneC^{-1-2\kappa}$, so $(\log N)^{-\gamma}\bigl(3B^{(1)}_N-3B^{(2)}_N\bigr)\to0$.
By the algebraic identity
\begin{equation}\label{eq:bookkeep}
3B^{(1)}_N-3B^{(2)}_N=-3{:}\qoneLin_N^2{:}\,\qoneIqc_N+3{:}\qoneLin_N^2{:}\,v_N+9c_{N,2}v_N,
\end{equation}
the singular products are exactly cancelled at the cost of the explicit
counterterm $+9c_{N,2}v_N$. Comparing with the counterterm $+9c_{N,2}u_N$ built
into the definition of $B_\gamma$, and using $u_N=\qoneLin_N-\qoneIqc_N+v_N$, the
\emph{residual} counterterm is
$9c_{N,2}u_N-9c_{N,2}v_N=9c_{N,2}(\qoneLin_N-\qoneIqc_N)$. This residual is multiplied by
$(\log N)^{-\gamma}$, and since $c_{N,2}\asymp\log N$ (Lemma~\ref{lem:logdiv}) it
contributes $(\log N)^{1-\gamma}\langle\qoneLin_N-\qoneIqc_N,\psi\rangle$, which does
\emph{not} vanish; hence the counterterm in $B_\gamma$ must be tuned to
$+9c_{N,2}v_N$ (equivalently $\qoneHp_3(\qonePN u)$ must be paired so that the residual
is absorbed) for Step~1 to close. This counterterm-matching is precisely the
content of issue~q1-2 and is outside the scope of the present lemma; granting that
matching, the singular contributions are handled by \eqref{eq:bookkeep} and
Lemma~\ref{lem:std}(ii), and pairing with $\psi$ gives
\[
\bigl\langle(\log N)^{-\gamma}\bigl(\qoneHp_3(u_N;c_N)+9c_{N,2}u_N\bigr),\psi\bigr\rangle\xrightarrow[N\to\infty]{}0\quad\text{a.s.},
\]
i.e.\ $u(t)\in B_\gamma$ for fixed $t$. Hence $\mu(B_\gamma)=1$.

\section*{Step 2: $(T_\psi^*\mu)(B_\gamma)=0$}

We must show $u+\psi\notin B_\gamma$. Expanding with $\psi_N=\qonePN\psi$,
\begin{align*}
&(\log N)^{-\gamma}\qoneHp_3(u_N+\psi_N;c_N)+(\log N)^{-\gamma}9c_{N,2}(u_N+\psi_N)\\
&\quad=\Bigl[(\log N)^{-\gamma}\qoneHp_3(u_N;c_N)+9(\log N)^{-\gamma}c_{N,2}u_N\Bigr]\\
&\qquad+3(\log N)^{-\gamma}{:}u_N^2{:}\,\psi_N+3(\log N)^{-\gamma}u_N\psi_N^2\\
&\qquad+(\log N)^{-\gamma}\psi_N^3+9(\log N)^{-\gamma}c_{N,2}\psi_N.
\end{align*}
The bracketed term tends to $0$ a.s.\ (Step 1). Since ${:}u_N^2{:}$ and $u_N$
converge to finite distributional limits (Lemma~\ref{lem:std}(i)), the next three
terms tend to $0$ after multiplication by $(\log N)^{-\gamma}$. The decisive term
is the last: by Lemma~\ref{lem:logdiv},
\[
9(\log N)^{-\gamma}c_{N,2}\,\langle\psi_N,\psi\rangle\ \gtrsim\ (\log N)^{-\gamma+1}\,\langle\psi_N,\psi\rangle .
\]
As $N\to\infty$, $\langle\psi_N,\psi\rangle\to\|\psi\|_{L^2}^2>0$ (because
$\psi\not\equiv0$), and $(\log N)^{-\gamma+1}\to\infty$ for $\gamma<1$. Choosing
$\tfrac12<\gamma<1$, the displayed quantity \emph{diverges}, so the full pairing
does not converge to $0$; that is, $u+\psi\notin B_\gamma$ a.s. Therefore
$(T_\psi^*\mu)(B_\gamma)=\mu(\{u:u+\psi\in B_\gamma\})=0$.

\section*{Conclusion}

By Steps 1 and 2, $\mu(B_\gamma)=1$ while $(T_\psi^*\mu)(B_\gamma)=0$. Thus $\mu$
and $T_\psi^*\mu$ are concentrated on disjoint Borel sets and are mutually
singular. The only structural inputs were that $\psi\not\equiv0$ (used solely to
ensure $\|\psi\|_{L^2}^2>0$) and that the coupling is nonzero, so the conclusion
holds for all such $\psi$. \qed

\section*{Remarks}

\begin{enumerate}\setlength\itemsep{4pt}
\item \textbf{Where dimension three enters.} The entire effect is carried by
$c_{N,2}\gtrsim\log N$ (Lemma~\ref{lem:logdiv}), the renormalization constant of
the resonance $\qoneLsq\circ\qoneIqc$. This logarithmic
divergence is absent in $d=2$, where $\mu$ \emph{is} quasi-invariant under smooth
shifts and even equivalent to the free field. The borderline dimension for
shift-quasi-invariance is exactly $3$, and the question is on which side it falls;
the answer is the singular side.

\item \textbf{Why the heuristic ``density'' argument fails.} Formally
$dT_\psi^*\mu/d\mu$ involves a term proportional to the would-be field
$\qoneLin^3-C\qoneLin$, which does \emph{not} define a random distribution in $d=3$ (its
covariance behaves like $|x-y|^{-3}$, non-integrable). The log-divergent constant
in the formal density is precisely what makes the candidate density blow up. The
rigorous proof above replaces the nonexistent density by the cut-off functional
defining $B_\gamma$ and extracts the divergence cleanly.

\item \textbf{Role of $\gamma$.} The window $\tfrac12<\gamma<1$ is forced:
$\gamma>\tfrac12$ kills the cube and the chaos remainders (Lemma~\ref{lem:cube}),
while $\gamma<1$ lets the shift term $(\log N)^{-\gamma}c_{N,2}\asymp(\log
N)^{1-\gamma}$ diverge. Any $\gamma$ in this range distinguishes the measures.

\item \textbf{The combinatorial constant $3$.} The coefficient $3$ multiplying
$c_{N,2}$ in Lemma~\ref{lem:std}(ii) is $\binom{3}{1}$: the number of ways of
choosing which of the three legs of $\qoneIqc$ (or of the cube $\qoneLcb$) survives the
single Wick contraction against $\qoneLsq$. The companion factor
$\binom{2}{1}=2$ (interchange of the two legs of $\qoneLsq$) is absorbed into the
definition \eqref{eq:cN2} of $c_{N,2}$. The second-moment computation in the proof
of Lemma~\ref{lem:std}(ii) shows $3$ is the \emph{unique} constant removing the
$\log N$ divergence.
\end{enumerate}

\begin{thebibliography}{9}
\bibitem{HP} M.~Hairer, J.~Peroni.
  \emph{Quasi shift invariance of $\Phi^4$ measures}. To appear.
\bibitem{HKN24} M.~Hairer, S.~Kusuoka, H.~Nagoji.
  \emph{Singularity of solutions to singular SPDEs}.
  Preprint, arXiv:2409.10037, 2024.
\bibitem{Glimm68} J.~Glimm.
  \emph{Boson fields with the $\Phi^4$ interaction in three dimensions}.
  Comm.\ Math.\ Phys.\ \textbf{10} (1968), 1--47.
\bibitem{HM18} M.~Hairer, K.~Matetski.
  \emph{Discretisations of rough stochastic PDEs}.
  Ann.\ Probab.\ \textbf{46} (2018), 1651--1709.
\bibitem{EW24} S.~Esquivel, H.~Weber.
  \emph{A priori bounds for the dynamic fractional $\Phi^4$ model on $\qoneT^3$}.
  Preprint, arXiv:2411.16536, 2024.
\bibitem{CC18} R.~Catellier, K.~Chouk.
  \emph{Paracontrolled distributions and the $3$-dimensional stochastic
  quantization equation}. Ann.\ Probab.\ \textbf{46} (2018), 2621--2679.
\bibitem{Hairer14} M.~Hairer.
  \emph{A theory of regularity structures}. Invent.\ Math.\ \textbf{198} (2014),
  269--504.
\end{thebibliography}

\endgroup
